% page 220 du fac-simile (image p-219 du scan)
% bloc 154.9 x 233.3 mm ; marge gauche 8.0 mm
\pgc{-12.700mm}{0.000mm}{
\l{212}
\l{}
\l{\sou{hepato}. (\sou{anat.}) Vicero situita, che la homo e la precipua}
\l{\cel{3}vertebrozi, an la parto supra di la latero rekta di la}
\l{\cel{3}abdomino, organo qua sekrecas bilo, e produktas materio}
\l{\cel{3}sukroza de qua lu richigas la sango. - DEFIS.}
\l{}
\l{hepatiko. (bot.) Genero de \textquotedbl{}ranunkulei\textquotedbl{} de qua un speco re-}
\l{\cel{3}komendesas kontre la hepato-morbi. - L. hepatica. EFIS.}
\l{}
\l{hepatito. I. (\sou{mineral.}) Lapido hepatea. - II. (\sou{patol.}) In-}
\l{\cel{3}flamuro di la hepato. - DEF.}
\l{}
\l{\sou{heptaedro}. - (\sou{geom.})\cel{2}Solido sep-edra. - DEP.}
\l{}
\l{\sou{heptagono.}\cel{1}- (\sou{geom.})\cel{2}Poligono qua havas sep lateri e sep}
\l{\cel{3}anguli. - DEF.}
\l{}
\l{\sou{heraldo}. (\sou{olim)}. Oficiro di qua la rolo konsistis ye portar}
\l{\cel{3}la mesaji (defii, sumni, milito-deklari, e c.), agar la}
\l{\cel{3}proklami, regulizar la festichevalierala, sorgar pri la}
\l{\cel{3}blazoni. - DEFIRS.}
\l{}
\l{\sou{herbo.}\cel{2}(\sou{bot.)}\cel{2}I. Planto di qua la stipo esas ne-ligna e}
\l{\cel{3}maxim-multa-kaze verda. - II. Vejetanti verda qui kovras}
\l{\cel{3}la prati e quin onu uzas por nutrar la bestiaro. - EFIS.}
\l{}
\l{\sou{herbario}. Kolektajo ek planti sikigita e konservata inter}
\l{\cel{3}paper-folii. - DEFIRS.}
\l{}
\l{\sou{herbivora}. Qua su nutras per herbo. - DEFIS.}
\l{}
\l{\sou{herboro}. (\sou{bot.)}\cel{2}Planto medicinala, sikigita o preparita}
\l{\cel{3}(ye la formo ye qua olu vendesas da herboristo). - DEFIS.}
\l{}
\l{\sou{herdo}.\cel{2}Parto petro-plakizita di la kameno-sulo, sur qua}
\l{\cel{3}onu produktas fairo. - DE.}
\l{}
\l{\sou{heredar}. (\sou{trans.}) (\sou{Yuro-cienco)}. Posedeskar ulo per sucedo}
\l{\cel{3}di persono (generale ek la gepatri o parenti) qua mortis.}
\l{\cel{3}- (\sou{anke metaf.})}
\l{}
\l{\sou{herezio}. (\sou{religio katolika}). Doktrino di ula dogmati Kristana}
\l{\cel{3}qui esas kontrea a la kredo katolika e, do, kondamnita da}
\l{\cel{3}la eklezio katolika. - DEFIRS.}
\l{}
\l{\sou{herisar}. (\sou{trans.}) I. (aludante homo od animalo).\cel{2}Igar rekta}
\l{\cel{3}e rigida (sua hari, sua pili, sua plumi). - II. Garnisar}
\l{\cel{3}(objekto) ye kozi pikiva qui impedas tushar lu. - FS.}
\l{}
\l{\sou{herisono}. (\sou{zool.})\cel{2}Mamifero insektivora, di qua la dorso-pelo}
\l{\cel{3}esas kovrita ye pikivi longa e rigida ek pili aglomerita}
\l{\cel{3}quin lu herisas per bulatreskar kande onu proximigas su a}
\l{\cel{3}lu.- L.\cel{1}\sou{erinaceus}. - FS.}
}
